\documentclass[11pt]{article}
\usepackage{amsmath, amssymb, booktabs, geometry, hyperref, algorithm, algpseudocode, graphicx}
\geometry{letterpaper, margin=1in}
\hypersetup{hypertexnames=false}

\title{COMPASS: Convex Optimization of Mediated Polygenic Architecture via SNP Scores}
\author{}
\date{}

\begin{document}

\maketitle

\begin{abstract}
We present COMPASS, a framework for estimating gene- and context-specific mediated heritability from genome-wide association study (GWAS) summary statistics and sparse variant--gene annotations. COMPASS augments the infinitesimal, gene-summed stratified LD score regression (S-LDSC) model with a nuclear-norm-regularized matrix of gene-specific deviations, while retaining an explicit non-mediated residual term. Penalties are selected by ten-fold cross-validation over global LD-connected components. We applied official BaselineLD-adjusted S-LDSC and COMPASS to seven European-ancestry GWAS using public activity-by-contact links, brain chromatin peaks linked to expressed genes, and predicted eQTL probabilities. S-LDSC recovered microglial enrichment in Alzheimer disease and neuronal enrichment in schizophrenia, bipolar disorder, major depression, amyotrophic lateral sclerosis, and anxiety. Gene-specific deviations improved held-out prediction for several H3K27ac-linked analyses, most strongly for bipolar disorder, Parkinson disease, and schizophrenia, but were completely shrunk for major depression and anxiety. High-confidence predicted-eQTL annotations produced resolved microglial and oligodendrocyte-precursor coefficients in Parkinson disease; only Parkinson disease retained gene-specific deviations. These results show that context-level enrichment is more stable than gene-level decomposition and that gene-level claims require annotation-matched, LD-aware validation.
\end{abstract}

\section{Model}
We start with the same setup as MESC. Let $\omega_j$ denote the marginal effect of single nucleotide polymorphism (SNP) $j \in \{1, \dots, J\}$ on a complex trait. We model $\omega_j$ as a causal cascade mediated through genes $k \in \{1, \dots, K\}$:
\begin{equation}
    \omega_j = \sum_{k=1}^K \beta_{jk} \alpha_k + \gamma_j
\end{equation}
where $\beta_{jk}$ is the true, unobserved regulatory effect of SNP $j$ on gene $k$, $\alpha_k$ is the causal effect of gene $k$ on the trait, and $\gamma_j \sim \mathcal{N}(0, \tau_\gamma)$ captures non-mediated (direct or pleiotropic) effects. 

Instead of relying on empirical eQTLs, we utilize $L$ non-negative annotations $a_{jkl} \ge 0$ representing the predicted regulatory potential of SNP $j$ on gene $k$ via biological mechanism $l \in \{1, \dots, L\}$. We place a linear variance prior on the regulatory effects:
\begin{equation}
    \beta_{jk} \sim \mathcal{N}(0, V_{jk}) \quad \text{where} \quad V_{jk} = \sum_{l=1}^L w_l a_{jkl}
\end{equation}
Here, $w_l \ge 0$ represents the unknown heritability contribution weight of mechanism $l$.

\subsection{The infinitesimal Gaussian assumption and its limitations}

A simple assumption would be infinitesimal effects, $\alpha_k \sim \mathcal{N}(0, \tau_\alpha)$. Given that $\beta_{jk}$ and $\alpha_k$ are independent and mean-zero, the variance of their product is the product of their variances:
\begin{equation}
    \text{Var}(\beta_{jk} \alpha_k) = \mathbb{E}[\beta_{jk}^2] \mathbb{E}[\alpha_k^2] = V_{jk} \tau_\alpha
\end{equation}
Assuming independence across genes, we sum over genes, substitute the linear variance model for $V_{jk}$, and commute the sums:
\begin{equation}
    \text{Var}(\omega_j) = \sum_{k=1}^K V_{jk} \tau_\alpha + \tau_\gamma = \tau_\alpha \sum_{k=1}^K \left( \sum_{l=1}^L w_l a_{jkl} \right) + \tau_\gamma = \tau_\alpha \sum_{l=1}^L w_l \left( \sum_{k=1}^K a_{jkl} \right) + \tau_\gamma = \sum_{l=1}^L (\tau_\alpha w_l) A_{jl} + \tau_\gamma,
\end{equation}
where $A_{jl} \equiv \sum_{k=1}^K a_{jkl}$. The gene-aggregated matrix $A$ is therefore simply a $J \times L$ SNP annotation matrix: each column $A_{\cdot l}$ can be used as an ordinary S-LDSC annotation with coefficient $\theta_l = \tau_\alpha w_l$. Thus this model is almost trivial, involving just summing over genes.

\subsection{Gene-specific resolution via conditional variance}

To achieve gene-specific resolution we cannot rely on the infinitesimal Gaussian assumptions above. Here we consider conditioning the variance of the marginal effect on the specific causal trait variance of each gene, $s_k = \alpha_k^2 \ge 0$,
\begin{equation}
    \text{Var}(\omega_j \mid \alpha) = \sum_{k=1}^K s_k V_{jk} + \tau_\gamma
\end{equation}
Substituting this conditional variance into the expected GWAS $\chi^2$ statistic yields,
\begin{align}
    \mathbb{E}[\chi_j^2 \mid \alpha] &= 1 + N \sum_{m=1}^J r_{jm}^2 \text{Var}(\omega_m \mid \alpha) \nonumber \\
    &= 1 + N \sum_{m=1}^J r_{jm}^2 \left( \sum_{k=1}^K s_k \sum_{l=1}^L w_l a_{mkl} + \tau_\gamma \right) \nonumber \\
    &= 1 + N \sum_{k=1}^K \sum_{l=1}^L (s_k w_l) \mathcal{T}_{jkl} + N \tau_\gamma l_j
\end{align}
where $r_{jm}$ is LD, $l_j = \sum_m r_{jm}^2$ is the standard non-mediated LD score, and the deterministic, pre-computable, sparse 3D annotation-gene LD Tensor,
\begin{equation}
    \mathcal{T}_{jkl} = \sum_{m=1}^J r_{jm}^2 a_{mkl}.
\end{equation}

Direct optimization of $s$ and $w$ has two pathologies:
\begin{enumerate}
    \item \textbf{Scale Non-Identifiability:} For any scalar $c > 0$, the transformation $s \to s/c$ and $w \to c w$ leaves the product $B = s w^T$ unchanged, rendering the individual scales of gene importance and annotation importance unidentifiable without arbitrary anchor priors. This is addressable, for example, by constraining the $L_1$ norm of $w$ (or $s$) to be 1, but this does not resolve the second pathology: 
    \item \textbf{Non-Convexity:} The bilinear optimization is non-convex, containing saddle points and local minima where optimization routines can get trapped.
\end{enumerate}


\subsection{Convex Nuclear Norm Relaxation}

Let $B_{kl} = s_k w_l$ define the $K \times L$ matrix of mediated heritability contributions. We relax the strict rank-1 constraint $B = s w^T$ by treating $B \in \mathbb{R}^{K \times L}$ as a general matrix parameter regularized by the nuclear norm $\|B\|_* = \sum_i \sigma_i(B)$, where $\sigma_i(B)$ are the singular values of $B$.

We formulate the inference as a weighted least-squares regression minimized over the space of non-negative matrices:
\begin{equation}
    \min_{B \in \mathbb{R}_{+}^{K \times L}, \; \tau_\gamma \ge 0} \quad \sum_{j=1}^J \eta_j \left( \chi_j^2 - 1 - N \sum_{k=1}^K \sum_{l=1}^L B_{kl} \mathcal{T}_{jkl} - N \tau_\gamma l_j \right)^2 + \lambda \|B\|_*
\end{equation}
where $\eta_j$ are regression weights. The production hierarchical COMPASS analyses use $\eta_j=1$ to avoid response-dependent weighting during cross-validation; official S-LDSC comparisons use LDSC's native iterative weights for LD redundancy and heteroskedasticity.

This relaxation yields three (potential) advantages:
\begin{itemize}
    \item \textbf{Convexity:} Because the squared-error loss is convex with respect to $B$ and the nuclear norm is convex, the objective is convex. Any converged local minimum is therefore globally optimal, although correlated annotations can make the optimizer non-unique.
    \item \textbf{Discovery of Regulatory Modules:} Rather than enforcing regulatory homogeneity (where every causal gene obeys the exact same regulatory weight vector $w$), a low-rank solution $\text{rank}(B) = r > 1$ decomposes into $r$ distinct mechanistic archetypes via Singular Value Decomposition (SVD):
    \begin{equation}
        B = \sum_{i=1}^r \sigma_i u_i v_i^T
    \end{equation}
    Here, the right singular vectors $v_i \in \mathbb{R}^L$ represent distinct biological pathways (e.g., promoter-driven vs. distal Hi-C loop enhancer-driven), while the left singular vectors $u_i \in \mathbb{R}^K$ cluster disease genes by their physical mechanism of action.
    \item \textbf{Computational Tractability:} Because the number of annotations $L$ is much smaller than the number of genes ($K \approx 20,000$), computing the reduced SVD during optimization takes $O(K L^2)$ time.
\end{itemize}

\section{Cell-Type Indexing}

The index $l$ need not correspond to a biochemical mechanism such as promoter activity, enhancer activity, or chromatin contact class. The same derivation applies if $l$ indexes cell types or cellular contexts. In particular, consider the model
\begin{equation}
    \beta_{jkl} \sim \mathcal{N}(0, w_l a_{jkl})
\end{equation}
where $l$ indexes cell types. We now model, 
\begin{equation}
    \omega_j = \sum_{k=1}^K \sum_{l=1}^L \beta_{jkl} \alpha_{kl} + \gamma_j
\end{equation}
where $\alpha_{kl}$ is the causal effect of gene $k$ in cell type $l$ on the trait. 
We then have,
\begin{equation}
    \text{Var}(\omega_j|\alpha) = \sum_{k=1}^K \sum_{l=1}^L \alpha_{kl}^2 w_l a_{jkl} + \tau_\gamma = \sum_{k=1}^K \sum_{l=1}^L B_{kl} a_{jkl} + \tau_\gamma
\end{equation}
where now $B_{kl} = \alpha_{kl}^2 w_l$ is the gene-by-cell-type mediated heritability matrix. The same convex nuclear norm relaxation applies, and the SVD of $B$ now identifies cell-type-specific regulatory modules.

\subsection{Hierarchical context and gene-deviation model}

Applying the nuclear-norm penalty directly to the complete matrix $B$ also penalizes a context-wide effect shared by all genes. This differs from the infinitesimal S-LDSC reduction, in which every gene contributing to annotation $l$ shares one unpenalized coefficient. To separate these two levels, we write
\begin{equation}
    B_{kl} = s\,q_l + D_{kl}, \qquad s\ge 0,\quad q_l\ge 0,\quad s q_l+D_{kl}\ge 0,
    \label{eq:hierarchical-B}
\end{equation}
where $q$ is a context profile, $s$ is its fitted genome-wide scale, and $D$ contains gene-specific departures from that profile. The resulting objective is
\begin{equation}
    \min_{s,\,D,\,\tau_\gamma\ge0:\;s q_l+D_{kl}\ge0}
    \sum_j \eta_j\left[
      \chi_j^2-1-N\sum_{k,l}(s q_l+D_{kl})\mathcal{T}_{jkl}
      -N\tau_\gamma l_j
    \right]^2 + \lambda\lVert D\rVert_*.
    \label{eq:hierarchical-objective}
\end{equation}
Thus, $D$ can either upweight or downweight a gene relative to its context-wide coefficient, while total gene-context heritability remains non-negative. Large $\lambda$ suppresses only gene-specific deviations and leaves the context-wide component estimable. The context design is the gene sum of the same annotation tensor, $C_{jl}=\sum_k a_{jkl}$, so that the fitted predictor is exactly $A(sq+D)$ rather than a sum of unmatched flat-context and gene-linked designs. When $q$ is unconstrained, the model estimates one non-negative coefficient per context. For the BaselineLD-calibrated analysis, $q$ is the vector of official S-LDSC coefficients after setting negative entries to zero; COMPASS estimates $s$ rather than transferring the reference-specific coefficient magnitude. Chromosome jackknife estimates are used for uncertainty in $s q_l$, with the S-LDSC profile uncertainty propagated approximately.

Jointly re-estimating $s$ at every value of $\lambda$ creates an avoidable allocation ambiguity: when the deviation penalty is weak, $D$ can reproduce broad context signal and drive $s$ to zero. We therefore use a sequential hierarchical estimator. Within each training fold, we first estimate $s$ and $\tau_\gamma$ at the context-only end of the path, where $D=0$. We then hold the resulting $s q$ fixed while cross-validating $\lambda$ for $D$; only $D$ and $\tau_\gamma$ are updated. The final full-data fit follows the same procedure. Consequently, the shared context effect has one definition across the regularization path, and held-out improvements measure gene-specific signal conditional on that shared component rather than reallocating it.

The proximal step begins at a common high-throughput value and uses objective-based backoff. A single uphill update is allowed because alternating the proximal and exact coordinate updates can produce a transient increase before descending. If the penalized objective remains above the best accepted value for five consecutive updates, that state is restored and the step is halved. This permits larger steps at strongly regularized points while controlling sustained instability at weaker penalties.

The calibrated analysis is deliberately a two-stage conditional model. Because $q$ is estimated from the same GWAS, its cell-type ordering is not an independent COMPASS discovery, and cross-validation selects only the penalty on $D$ conditional on $q$. The analysis asks whether gene-specific departures improve held-out prediction while retaining a BaselineLD-adjusted context contrast; it is not an unbiased end-to-end validation of that contrast.

For interpretation, we report the context-wide and gene-specific contributions separately. If $M_l=\sum_{j,k}a_{jkl}$ is the gene-linked annotation mass in context $l$, the direct heritability attributed to context $l$ is
\begin{equation}
    H_l^{\mathrm{global}}=M_l s q_l,\qquad
    H_l^{\mathrm{deviation}}=\sum_{j,k}a_{jkl}D_{kl},\qquad
    H_l^{\mathrm{total}}=H_l^{\mathrm{global}}+H_l^{\mathrm{deviation}}.
\end{equation}
This decomposition tests whether a weakly penalized deviation matrix changes the aggregate cell-type ranking even though the global component is fixed.

\subsection{The Necessity of Non-Negativity}
While standard summary statistic methods occasionally tolerate negative enrichment parameters, enforcing $B_{kl} \ge 0$ is mandatory in this gene-centric framework. Without non-negativity constraints, the optimization landscape suffers from two failures:

\paragraph{LD Collinearity Overfitting.}
Neighboring genes on a chromosome reside within the same LD block, causing their annotation tensors to be highly collinear ($\mathcal{T}_{j, k_1, l} \approx \mathcal{T}_{j, k_2, l}$ for adjacent genes $k_1, k_2$). If an unconstrained model attempts to fit a localized GWAS signal corresponding to a true heritability contribution of $+10$, it can overfit statistical noise by assigning extreme, mutually cancelling coefficients to adjacent genes, e.g. $B_{k_1, l} = +100$ and $B_{k_2, l} = -90$ so that $B_{k_1, l} + B_{k_2, l} = +10$. Such cancellation preserves the aggregate fit while making gene-level coefficients unstable. The constraint $B_{kl} \ge 0$ prevents this signed cancellation, although it cannot by itself identify one gene among collinear alternatives.

\paragraph{Mathematical and Biological Plausibility}
By definition, $B_{kl} = \alpha_k^2 w_l$ represents a variance contribution mediated by gene $k$ via mechanism $l$ and therefore cannot be negative. Unconstrained negative entries can also offset positive LD-tagged variance and may produce fitted expectations below the null value. The hierarchical model permits negative deviations only relative to the shared context effect and retains the required constraint on the total, $s q_l+D_{kl}\ge0$.

\subsection{Optimization via Alternating Projections}

Because standard singular value thresholding (SVT) does not preserve elementwise lower bounds, the proximal operator for the nuclear norm plus the feasible box is evaluated by Dykstra's alternating proximal algorithm.

Let $Z$ denote either the nuclear-model matrix $B$, with lower bound $L=0$, or the hierarchical deviation matrix $D$, with lower bound $L_{kl}=-s q_l$. Let $f(Z,\tau_\gamma)$ denote the squared-error loss. One proximal-gradient update proceeds as follows:

\begin{algorithm}[H]
\caption{Box-constrained nuclear-norm proximal update}
\begin{algorithmic}[1]
\State \textbf{Input:} Current $Z^{(t)}$, lower-bound matrix $L$, step size $\zeta$, penalty $\lambda$
\State $G \gets Z^{(t)}-\zeta\nabla_Z f(Z^{(t)},\tau_\gamma^{(t)})$
\State $X\gets G$; $P\gets0$; $Q\gets0$
\Repeat
    \State $Y\gets\operatorname{SVT}_{\zeta\lambda}(X+P)$; $P\gets X+P-Y$
    \State $X_{\mathrm{new}}\gets\max(L,Y+Q)$; $Q\gets Y+Q-X_{\mathrm{new}}$
    \State $X\gets X_{\mathrm{new}}$
\Until{the Dykstra residual converges}
\State $Z^{(t+1)}\gets X$
\State Update $\tau_\gamma^{(t+1)}$ by exact non-negative coordinate minimization
\State \textbf{Output:} feasible low-rank update $Z^{(t+1)}$ and residual coefficient $\tau_\gamma^{(t+1)}$
\end{algorithmic}
\end{algorithm}

For the non-negative nuclear model, the implementation uses Dykstra iterations to evaluate the proximal operator of the nuclear norm plus the non-negative orthant. For the hierarchical model, the same procedure is applied to $D$ with the shifted box constraint $D_{kl}\ge -s q_l$, so negative deviations are allowed while total coefficients remain non-negative. The residual coefficient has an exact non-negative coordinate update, and objective-based step backoff controls sustained uphill trajectories. These operations solve a convex objective, but collinear annotations can still yield multiple equivalent decompositions; cross-validation and sensitivity analyses are therefore required for interpretation.

\subsection{Cross-validation strategy}

Leaving out an entire chromosome is not a valid penalty-selection scheme for determining $\lambda$. For example, if chromosome 1 is held out, then genes on chromosome 1 have no training annotations. Prediction on the held-out chromosome therefore relies on gene effects that were not estimable in the training set.

We therefore use the global, LD-aware procedure in Algorithm~\ref{alg:cre-cv}. The essential change is to define folds for genomic regions before considering genes. We construct a graph on variants, joining variants whose reference-panel LD exceeds a threshold $r_{jm}^2\ge\rho_{\mathrm{CV}}$, and use its connected components as indivisible holdout units. Each component receives one genome-wide fold label, so the same fold has the same meaning for every gene and a variant linked to multiple genes cannot receive conflicting labels. Components are distributed across folds to balance the number of regression rows and, where possible, to place different components linked to the same gene in different folds.

Let $\mathcal{C}$ denote the set of CRE--gene links, and write $j\in c$ when variant $j$ lies in CRE $c$. After the LD components have been assigned labels $z_q$, define
\begin{equation}
\begin{aligned}
\mathcal{K}(j)
    &= \{k:\exists\,(c,k)\in\mathcal{C}\text{ with }j\in c\}, \\
\mathcal{K}_{-p}
    &= \{k:\exists\,(c,k)\in\mathcal{C},\ \exists\,q\text{ with }c\cap D_q\ne\varnothing\text{ and }z_q\ne p\}, \\
\mathcal{S}_p
    &= \{j\in V_p:\mathcal{K}(j)\cap\mathcal{K}_{-p}\ne\varnothing\}.
\end{aligned}
\label{eq:cv-sets}
\end{equation}
Here, $\mathcal{K}(j)$ is the set of genes linked to variant $j$; $\mathcal{K}_{-p}$ is the set of ``training'' genes for fold $p$, i.e. that have at least one CRE overlapping a training LD component; and $\mathcal{S}_p$ is the set of held-out variants linked to at least one training gene. The intersection criterion allows a scored variant to have additional linked genes without training support. Those genes do not disqualify the variant: their effects remain unmodeled validation signal, making the score conservative rather than removing an otherwise informative row.

For fold $p$, all regression rows in $V_p$ are removed from the fitting loss. When scoring $j\in\mathcal{S}_p$, prediction is restricted to the supported genes,
\begin{equation}
\widehat{\chi}_{j,p\lambda}^{2,(-p)}
=1+N\sum_{k\in\mathcal{K}_{-p}}\sum_{l=1}^L\widehat{B}_{kl}^{(p,\lambda)}\mathcal{T}_{jkl}
+N\widehat{\tau}_{\gamma}^{(p,\lambda)}l_j.
\end{equation}
Variants outside CREs also inherit their component's fold and are held out of the training loss, preventing them from carrying correlated GWAS signal across the fold boundary, but they do not themselves contribute to the validation score used to choose $\lambda$.

\begin{algorithm}[H]
\caption{Global LD-component cross-validation for selecting $\lambda$}
\label{alg:cre-cv}
\begin{algorithmic}[1]
\State \textbf{Input:} CRE--gene links; GWAS variants; $A$ and $R^2$; descending penalty grid $\Lambda$; $F$ folds; LD threshold $\rho_{\mathrm{CV}}$
\State Construct $G=(V,E)$ with $V=\{1,\ldots,J\}$ and $E=\{(j,m):r_{jm}^2\ge\rho_{\mathrm{CV}}\}$
\State Let $D_1,\ldots,D_Q$ be the connected components of $G$
\State Assign each $D_q$ one global label $z_q\in\{1,\ldots,F\}$, balancing rows and spreading components linked to the same gene across folds
\State $\mathcal{P}\gets\varnothing$
\For{$p=1,\ldots,F$}
    \State $V_p\gets\bigcup_{q:z_q=p}D_q$ \Comment{All rows in held-out LD components}
    \State Compute $\mathcal{K}_{-p}$ and $\mathcal{S}_p$ from Equation~\ref{eq:cv-sets}
    \If{$\mathcal{S}_p=\varnothing$} \State \textbf{continue} \EndIf
    \State $\mathcal{P}\gets\mathcal{P}\cup\{p\}$
    \State Set regression weight $\eta_j^{(p)}\gets\eta_j\,\mathbf{1}[j\notin V_p]$
    \State Initialize $B\gets 0$ and $\tau_\gamma\gets 10^{-8}$
    \For{$\lambda\in\Lambda$ from largest to smallest}
        \State Fit $(B,\tau_\gamma)$ using weights $\eta^{(p)}$, warm-starting from the preceding $\lambda$
        \State Predict $\widehat{\chi}_{j,p\lambda}^{2,(-p)}$ using only genes in $\mathcal{K}_{-p}$
        \State $E_{p\lambda}\gets |\mathcal{S}_p|^{-1}\sum_{j\in\mathcal{S}_p}(\chi_j^2-\widehat{\chi}_{j,p\lambda}^{2,(-p)})^2$
    \EndFor
\EndFor
\State $\lambda^*\gets\arg\min_{\lambda\in\Lambda}|\mathcal{P}|^{-1}\sum_{p\in\mathcal{P}} E_{p\lambda}$
\State Refit on all regression rows, again warm-starting from larger penalties down to $\lambda^*$
\State \textbf{Output:} $\lambda^*$ and the full-data estimates $(\widehat B,\widehat\tau_\gamma)$
\end{algorithmic}
\end{algorithm}

This design directly targets the desired generalization: the model learns $B_{kl}$ from some regulatory regions for gene $k$ and tests whether those estimates predict summary statistics in other regions for that gene, while preventing pairs with $r_{jm}^2\ge\rho_{\mathrm{CV}}$ from being split across training and validation. Unlike chromosome holdout, the size of each holdout unit is determined by the observed LD graph rather than chromosome boundaries; unlike per-gene fold labeling, its label is globally defined.

\section{GWAS, annotations, and analysis design}

\subsection{GWAS panel}

The current analysis panel contains seven European-ancestry GWAS (Table~\ref{tab:gwas}). Every retained summary-statistic row has a known study-provided or design-derived effective sample size. GRCh38 releases were lifted to hg19 before matching to the LD references. The 2026 Alzheimer disease (AD) analysis supersedes the historical Jansen et al. 2019 analysis; the latter is retained only for method development and simulation diagnostics.

\begin{table}[ht]
\centering
\small
\caption{GWAS used in the current cross-trait analyses.}
\label{tab:gwas}
\resizebox{\textwidth}{!}{%
\begin{tabular}{lll}
\toprule
Trait & Release & Analysis note \\
\midrule
AD & 2026 consensus, GCST90704647 & European, proxy samples excluded; $N=1{,}119{,}443$ \\
Schizophrenia (SCZ) & Bigdeli et al. 2026 & European meta-analysis; variant-specific $N_{\mathrm{eff}}$ \\
Parkinson disease (PD) & Leonard et al. 2025 & European clinical-only; 34,933 cases, 31,009 controls \\
Bipolar disorder (BD) & PGC4, 2025 & European, excluding 23andMe \\
Major depression (MDD) & PGC, 2025 & European, excluding 23andMe \\
Amyotrophic lateral sclerosis (ALS) & Project MinE, 2021 & European; GCST90027164 \\
Anxiety & Friligkou et al. 2024 & European genomic-factor GWAS; $N=1{,}096{,}458$ \\
\bottomrule
\end{tabular}}
\end{table}

\subsection{Annotation families}

We evaluated four annotation families. First, the public Nasser et al. activity-by-contact (ABC) resource supplied CRE--gene links with ABC score at least 0.015. Brain and myeloid contexts were analyzed with and without adipose, skeletal-muscle, and uterine negative controls. Second, Glass Lab v2 supplied matched astrocyte, microglia, neuron, and oligodendrocyte ABC predictions. We analyzed continuous scores, binary support, and distal links after excluding self-promoters. Third, binary ATAC, H3K27ac, and H3K4me3 peaks from LHX2 astrocytes, PU1 microglia, NeuN neurons, and Olig2 oligodendrocytes were used directly for S-LDSC. For COMPASS, a variant in a peak was linked to each gene expressed in that cell type whose transcription start site was within 100 kb. Fourth, predicted eQTL probabilities from CatBoost/GPN models were available for astrocytes, excitatory neurons, inhibitory neurons, microglia, oligodendrocyte precursor cells (OPCs), and oligodendrocytes. The primary predicted-eQTL analysis retained links with probability at least 0.9 and used their probabilities as annotation values.

All COMPASS fits retain every GWAS variant represented in the UK Biobank LD reference, including variants with no annotation. Such variants have all-zero mediated annotations but remain in the regression and in the non-mediated residual LD-score term. UK Biobank LD is represented by chromosome-level sparse $r^2$ matrices: entries below $r^2=0.01$ are set to zero, retained values are stored in float16, and model coefficients are optimized in float32. The $r^2\ge0.01$ AD graph contained 3,354 connected components across 12.36 million variants; ten folds were balanced to within one row.

The equations use a common $N$ for clarity. In every real-data fit, the corresponding row is scaled by its known variant-specific $N_j$ when available; fixed effective sample sizes are used only for releases that do not report per-variant values.

Official S-LDSC analyses use the 1000 Genomes Phase 3 European reference, BaselineLD v2.2, HapMap3 regression SNPs, matching allele frequencies, a 1 cM LD window, overlap-annotation accounting, native iterative weights, and block-jackknife standard errors. Continuous gene-linked annotations are summed over genes before LD-score calculation. The primary context statistic is the coefficient divided by its jackknife standard error; enrichment is reported as a secondary descriptive quantity because overlapping annotations can have significant marginal enrichment despite an unresolved conditional coefficient.

For hierarchical COMPASS, negative S-LDSC coefficients are set to zero to define $q$, COMPASS re-estimates a single non-negative scale $s$, and signed gene-specific deviations $D$ are fit subject to $s q+D\ge0$. The S-LDSC profile and COMPASS deviations use the same cell labels, but the brain-peak profile is estimated from flat H3K27ac peaks whereas the COMPASS tensor links those peaks to nearby expressed genes. Thus context ordering is supplied by S-LDSC and gene-level prediction is tested by COMPASS; this is a two-stage conditional analysis, not independent replication of the context profile. The predicted-eQTL analysis uses the gene sum of the same probabilistic tensor in both stages.

\section{Results}

\subsection{ABC analyses and model diagnostics}

The initial Jansen 2019 AD analysis used the public ABC panel to evaluate the nuclear relaxation. Ten-fold LD-component CV selected $\lambda=100$ (mean MSE 2.98059 versus 2.98074 at $\lambda\ge300$). The fitted matrix was numerically rank one, with leading singular value $8.63\times10^{-7}$ and all subsequent values near $10^{-14}$. A strict rank-one fit initialized from the nuclear solution had coefficient correlation 0.99998 and relative Frobenius difference 0.0139. Adding adipose, muscle, and uterus controls again selected $\lambda=100$ and a rank-one solution. Controls received 8.6--10.6\% of coefficient mass, comparable to the 6.9--11.1\% assigned to brain and myeloid contexts. This failure of the intended controls demonstrates that the original nuclear model did not identify cell-type-specific ABC effects.

Agora cross-reference and pathway analyses of these historical fits were consistent with broad rather than specific prioritization. Among 18,396 assayed genes, 865 overlapped 967 Agora nominated targets. At the top 1,000 genes, the nuclear and rank-one fits contained 68 and 69 Agora targets (approximately 47 expected; unadjusted $p=0.00139$ and $0.000877$), whereas both top-100 lists contained six targets ($p=0.33$). Context rankings were nearly identical and g:Profiler terms were dominated by broad tissue signatures. These comparisons support weak aggregate overlap but not validated AD gene ranking.

The historical official-LDSC diagnostic used the same 12.36 million UK Biobank variants, gene-summed ABC scores, and an all-SNP baseline. It estimated observed-scale $h^2=0.016$ (SE 0.002), intercept 1.0155 (SE 0.0062), bipolar-neuron annotated $h^2=0.0028$ (SE 0.0012; coefficient $z\simeq2.3$), and CD14 annotated $h^2=0.0015$ (SE 0.0010). Under the conventional 1000 Genomes/BaselineLD analysis, the bipolar-neuron coefficient fell to $z=0.60$ and CD14 to $z=-0.08$. The apparent matched-reference bipolar signal was therefore not robust to standard baseline adjustment and is retained only as a reference-panel diagnostic.

In the 2026 GWAS, BaselineLD-adjusted public ABC coefficients were unresolved for AD (largest positive $z=1.34$ for bipolar-neuron ABC) and for SCZ. Glass continuous and distal ABC scores were also unresolved: the largest positive conditional $z$ was 1.56 for distal microglial AD links and 1.31 for continuous oligodendrocyte SCZ links. Binarization was more informative. Binary Glass microglial support had $z=3.22$ and 27.1-fold enrichment in AD, while binary Glass neuronal support had $z=3.74$ and 10.6-fold enrichment in SCZ. Self-promoters contributed 86--87\% of continuous Glass ABC mass, explaining why continuous scores were largely redundant with BaselineLD promoter features. Distal support was cell-specific (pairwise Jaccard indices approximately 0.001) but too sparse for resolved conditional coefficients.

Uncalibrated nuclear COMPASS did not recover the binary S-LDSC specificity. Public-ABC fits selected $\lambda=30$ for AD and $\lambda=3$ for SCZ, but adding peripheral controls left each control with 8--11\% of coefficient mass. Continuous Glass fits assigned approximately equal mass to all four cell types (AD 21--22\%; SCZ 19--21\%). Peak-to-expressed-TSS nuclear fits were similarly diffuse. These results motivated the hierarchical decomposition and prevent interpreting the uncalibrated nuclear context weights as cell-type enrichments.

\subsection{Sparse ABC simulation}

We simulated $h^2=0.20$ using real ABC and LD, assigning 70\% to bipolar-neuron annotations, 30\% to CD14-monocyte annotations, and zero to adipose, muscle, and uterus. Only 10\% of ABC-positive variants in each causal context had direct effects. Three replicates were generated at $N_{\mathrm{eff}}=100{,}000$ and $200{,}000$. The nuclear model ranked both causal contexts above all controls in two of three lower-$N$ replicates and all three higher-$N$ replicates, but assigned only 46--47\% of normalized fitted mass to the two causal contexts; null controls retained 53--54\%. The mean $L_1$ error for the causal 70:30 split was approximately 0.54. In contrast, the joint infinitesimal S-LDSC baseline assigned positive coefficients to both causal annotations and near-zero or negative coefficients to controls, although its total heritability and residual term were miscalibrated under the deliberately sparse generative model. The simulation therefore identifies annotation and LD collinearity, rather than optimization alone, as a major limit on cell-type recovery.

\subsection{Brain chromatin enrichments across seven GWAS}

Table~\ref{tab:h3-sldsc} summarizes the joint H3K27ac coefficients after BaselineLD adjustment. AD was microglia-predominant: PU1 H3K27ac had $z=2.97$, 13.0-fold enrichment, and enrichment $p=2.15\times10^{-7}$. SCZ and BD showed strong neuronal coefficients ($z=7.28$ and 7.74), as did MDD ($z=6.68$) and anxiety ($z=4.77$). ALS had a smaller neuronal coefficient ($z=2.49$). The clinical-only PD coefficients were individually unresolved, with neuronal $z=1.68$ and microglial $z=0.99$. MDD also had an astrocyte coefficient ($z=2.67$). The complete 12-annotation fits included ATAC and H3K4me3; notable concordant results included NeuN H3K4me3 in SCZ ($z=4.63$) and BD ($z=4.17$), whereas no ATAC coefficient added a comparably strong independent signal.

\begin{table}[ht]
\centering
\small
\caption{Conditional coefficient $z$-scores for four H3K27ac annotations in joint BaselineLD-adjusted S-LDSC. A, astrocyte (LHX2); N, neuron (NeuN); M, microglia (PU1); O, oligodendrocyte (Olig2).}
\label{tab:h3-sldsc}
\begin{tabular}{lrrrr}
\toprule
GWAS & A & N & M & O \\
\midrule
AD 2026 & $-0.14$ & $0.61$ & $2.97$ & $0.62$ \\
SCZ 2026 & $1.53$ & $7.28$ & $-1.03$ & $1.33$ \\
PD 2025 & $-0.13$ & $1.68$ & $0.99$ & $0.16$ \\
BD 2025 & $1.52$ & $7.74$ & $-0.02$ & $-0.02$ \\
MDD 2025 & $2.67$ & $6.68$ & $0.86$ & $1.93$ \\
ALS 2021 & $-0.43$ & $2.49$ & $0.16$ & $1.11$ \\
Anxiety 2024 & $-0.54$ & $4.77$ & $0.49$ & $0.40$ \\
\bottomrule
\end{tabular}
\end{table}

\subsection{Hierarchical COMPASS with H3K27ac-linked genes}

The corrected signed hierarchical model selected nonzero gene-specific deviations for AD, SCZ, PD, BD, and ALS (Table~\ref{tab:h3-compass}). Held-out improvement over the context-only limit ranged from 0.82\% for ALS to 5.01\% for BD. MDD and anxiety selected the context-only solution and had exactly zero deviations. The selected AD result supersedes an earlier one-sided, flat-context pilot that had selected the no-deviation limit. In the corrected annotation-matched design, AD selected $\lambda=100$, improved mean MSE by 0.89\%, and remained microglia-predominant (53.3\% of total modeled contribution). SCZ selected $\lambda=100$, improved MSE by 2.41\%, and assigned 50.2\% to neurons. PD and BD had the largest relative improvements, 4.73\% and 5.01\%, respectively. PD divided most modeled mass between neurons (52.0\%) and microglia (42.3\%); BD assigned 38.9\% to neurons and 32.3\% to the gene-shared intercept. ALS assigned 39.0\% to neurons, 26.7\% to the intercept, and 19.3\% to oligodendrocytes.

\begin{table}[ht]
\centering
\scriptsize
\caption{Corrected signed hierarchical COMPASS results for H3K27ac peak-to-expressed-gene annotations. MSE$_0$ is the context-only limit. Fractions include the gene-shared intercept (I).}
\label{tab:h3-compass}
\resizebox{\textwidth}{!}{%
\begin{tabular}{lrrrrll}
\toprule
GWAS & $\lambda^*$ & MSE & MSE$_0$ & Improvement & Largest total contributions & Deviation \\
\midrule
AD & $100$ & 16.4093 & 16.5573 & 0.89\% & M 53.3\%, I 19.2\%, N 11.0\% & rank 2 \\
SCZ & $100$ & 13.8820 & 14.2247 & 2.41\% & N 50.2\%, I 21.5\%, A 11.6\% & rank 2 \\
PD & $100$ & 12.9076 & 13.5487 & 4.73\% & N 52.0\%, M 42.3\%, O 5.7\% & rank 1 \\
BD & $100$ & 7.4173 & 7.8086 & 5.01\% & N 38.9\%, I 32.3\%, A 12.5\% & rank 2 \\
MDD & $10^5$ & 9.3957 & 9.3957 & 0.00\% & N 61.0\%, A 19.0\%, O 13.9\% & zero \\
ALS & $100$ & 2.8897 & 2.9137 & 0.82\% & N 39.0\%, I 26.7\%, O 19.3\% & rank 2 \\
Anxiety & $10^6$ & 6.3371 & 6.3371 & 0.00\% & N 85.6\%, M 8.2\%, O 6.2\% & zero \\
\bottomrule
\end{tabular}}
\end{table}

Signed deviations were common in the nonzero fits: AD contained 31,238 positive and 29,873 negative entries; SCZ contained 30,058 positive and 29,989 negative entries. Every total coefficient satisfied $s q+D\ge0$. A matched SCZ ablation constrained $D\ge0$ and also selected $\lambda=100$ (MSE 13.9287). The signed fit had 0.34\% lower MSE, but the paired fold difference was unresolved (paired $t$-test $p=0.667$). The signed model nevertheless had lower full-data objective (9.2668 versus 9.7025) and two resolved singular values rather than one. Thus the signed formulation is biologically appropriate and fits the data better, but its predictive advantage over non-negative deviations is not established by this single comparison.

Contribution-weighted ordered g:Profiler analyses used the smallest leading gene sets accounting for 25\% and 50\% of positive deviation contribution, with all assayed genes as background. AD global and microglial rankings emphasized osteoclast differentiation, Fc receptor and B-cell signaling, lipid/cholesterol processes, and BDNF/EGFR signaling. SCZ neuronal and global rankings emphasized synapse organization, trans-synaptic signaling, neurogenesis, neuronal systems, Rett syndrome, and 16p11.2 deletion, while the microglial ranking emphasized cytokine and interferon signaling. PD rankings included MAPK, IL2, VEGF, and lysosome-related terms. BD rankings included cognition, learning or memory, and fatty-acid metabolism. ALS rankings were dominated by broad neuronal, glial, and tissue-expression terms. MDD and anxiety had no deviation-ranked genes. Because each trait produces multiple overlapping ordered queries and no external replication set was used, these g:SCS-corrected terms are exploratory descriptions of fitted rankings rather than additional association tests.

\subsection{High-confidence predicted-eQTL annotations}

Predicted-eQTL links were analyzed for PD, BD, MDD, ALS, and anxiety. Source GRCh38 positions were lifted once to hg19, and failed, cross-chromosome, or multiply mapped positions were discarded. At probability at least 0.9, 12,882,024 source links passed the threshold and 11,936,542 unique lifted links were retained before GWAS matching. Table~\ref{tab:eqtl-sldsc} reports conditional S-LDSC coefficients. PD showed positive microglial ($z=2.92$) and OPC ($z=2.68$) coefficients. No coefficient reached $|z|=2$ in the other traits; the largest were excitatory-neuron BD ($z=1.75$), inhibitory-neuron MDD ($z=-1.70$), microglial ALS ($z=1.50$), and microglial anxiety ($z=1.37$). These conditional results are more conservative than the corresponding marginal enrichment statistics, which were positive for many correlated annotations.

\begin{table}[ht]
\centering
\scriptsize
\caption{Conditional coefficient $z$-scores for continuous predicted-eQTL annotations in joint BaselineLD-adjusted S-LDSC. Exc, excitatory neuron; Inh, inhibitory neuron.}
\label{tab:eqtl-sldsc}
\begin{tabular}{lrrrrrr}
\toprule
GWAS & Astrocyte & Exc & Inh & Microglia & OPC & Oligodendrocyte \\
\midrule
PD & $0.49$ & $-1.41$ & $-0.15$ & $2.92$ & $2.68$ & $-0.17$ \\
BD & $0.35$ & $1.75$ & $0.93$ & $0.20$ & $0.98$ & $-0.70$ \\
MDD & $1.24$ & $0.87$ & $-1.70$ & $0.52$ & $1.35$ & $-0.56$ \\
ALS & $-0.46$ & $0.06$ & $0.96$ & $1.50$ & $-0.46$ & $1.32$ \\
Anxiety & $-0.64$ & $0.22$ & $1.29$ & $1.37$ & $-0.41$ & $-0.12$ \\
\bottomrule
\end{tabular}
\end{table}

Only PD retained gene-specific predicted-eQTL deviations (Table~\ref{tab:eqtl-compass}). It selected $\lambda=100$, lowering mean held-out MSE from 133.882 to 132.836 (0.78\%). OPCs accounted for 50.5\%, microglia 34.1\%, and astrocytes 15.4\% of modeled context contribution. The 25\% and 50\% ranked gene sets yielded 12 and 21 g:SCS-significant terms, including global synaptic-vesicle endocytosis, recycling, and pathway terms; these are exploratory because the same GWAS selected and ranked the deviations. BD, MDD, and ALS selected $\lambda=10^6$ and exactly zero deviations. Anxiety nominally selected $\lambda=30$, but its score differed from the high-penalty plateau only in the sixth decimal place and the final deviation matrix was zero. Their reported context fractions therefore reproduce the clipped S-LDSC profile after annotation-mass scaling and do not constitute COMPASS gene-level discoveries.

\begin{table}[ht]
\centering
\small
\caption{Hierarchical COMPASS results for predicted-eQTL probability annotations.}
\label{tab:eqtl-compass}
\begin{tabular}{lrrrrl}
\toprule
GWAS & $\lambda^*$ & MSE & MSE$_0$ & Improvement & Largest total contributions \\
\midrule
PD & $100$ & 132.836 & 133.882 & 0.78\% & OPC 50.5\%, M 34.1\%, A 15.4\% \\
BD & $10^6$ & 8.3650 & 8.3650 & 0.00\% & Exc 39.4\%, Inh 28.7\%, OPC 18.1\% \\
MDD & $10^6$ & 15.1278 & 15.1278 & 0.00\% & A 44.4\%, OPC 26.5\%, Exc 24.0\% \\
ALS & $10^6$ & 3.8401 & 3.8401 & 0.00\% & O 41.4\%, Inh 35.0\%, M 22.1\% \\
Anxiety & $30$ & 7.39170 & 7.39170 & $<0.01$\% & Inh 64.2\%, M 25.7\%, Exc 10.1\% \\
\bottomrule
\end{tabular}
\end{table}

\section{Discussion}

Across annotation families and GWAS, context-level S-LDSC signals were more stable than gene-level decompositions. Brain H3K27ac annotations recovered expected disease biology: microglia for AD and neurons for SCZ, BD, MDD, ALS, and anxiety. Predicted eQTL probabilities added a distinct PD signal in microglia and OPCs. By contrast, continuous ABC scores were usually unresolved after BaselineLD adjustment, while binary ABC support reproduced AD microglial and SCZ neuronal enrichment. This difference is consistent with promoter-dominated continuous ABC mass and indicates that annotation scaling and support can matter as much as the underlying CRE--gene map.

Hierarchical COMPASS found held-out gene-specific signal for five of seven H3K27ac-linked analyses, but effect sizes were modest and only PD retained deviations under predicted-eQTL annotations. The two-stage design deliberately conditions on an S-LDSC-derived context profile estimated from the same GWAS. Consequently, the context fractions are not independent COMPASS discoveries; the independent quantity tested by CV is improvement from $D$. Nuclear regularization creates low-rank rather than elementwise-sparse solutions, so thousands of positive and negative coefficients can contribute even when only a few latent factors are resolved. Pathway results therefore depend on contribution-weighted truncation and require external replication.

Several safeguards are essential. All regression variants, not only annotated variants, must be retained to model LD and the residual term. Whole-chromosome CV is invalid because genes on the held-out chromosome have no training annotations; global LD-component folds instead separate correlated variants while allowing different regulatory regions for the same gene to appear in training and validation. The $r^2=0.01$ sparsification threshold makes chromosome-level LD tractable, but correlations below that threshold are omitted. Finally, S-LDSC and COMPASS use different reference panels and weighting schemes, and the brain-peak S-LDSC profile is estimated from flat peaks before COMPASS links those peaks to genes. These differences make the combined estimator a calibrated conditional workflow rather than one joint likelihood.

The current results support cell-type enrichment for several traits and gene-level predictive structure in selected annotation--GWAS combinations, but they do not support a universal gene-prioritization model. Stronger validation will require independent GWAS, independent regulatory maps, or simulations whose causal variant--gene effects are generated outside the fitted annotation family.

\end{document}
